\documentclass[letter, 11pt]{article}
\usepackage[utf8]{inputenc}
\usepackage[english]{babel}
\usepackage{amsfonts}
\usepackage{amsmath}
\usepackage[dvips]{graphicx}
\usepackage{url}
\usepackage[top=3cm,bottom=3cm,left=3.5cm,right=3.5cm,footskip=1.5cm,headheight=1.5cm,headsep=.5cm,textheight=3cm]{geometry}
\usepackage{listings}
\usepackage{color}
\usepackage{fancyvrb}
\usepackage{fancyhdr}
\usepackage{enumerate}
\usepackage{amssymb}
\usepackage{multirow}

%%%%%%%%%%%%%%%%%%%%%%
%Estilo del documento%
%%%%%%%%%%%%%%%%%%%%%%
\pagestyle{fancyplain}

%%%%%%%%%%%%%%%%%%%%%%%%%%%%%%%%%%%%%%%%%%%
%Fancyheadings. Top y Bottom del documento%
%%%%%%%%%%%%%%%%%%%%%%%%%%%%%%%%%%%%%%%%%%%
% Recuerde que en este documento la portada del documento no posee
% numeracion, pero de igual manera llamaremos a esa primera pagina la numero
% 1, y la que viene la dos. Esto es para tener una idea de las que
% llamaremos pares e impares
\lhead{} %Parte superior izquierda
\rhead{\bf \it UTFSM, Casa Central} %Parte superior derecha
\lfoot{JSCO} %Parte inferior izquierda.
\cfoot{} %Parte inferior central
\rfoot{\bf \it Page \thepage} %Parte inferior derecha
\renewcommand{\footrulewidth}{0.4pt} %Linea de separacion inferior


\begin{document}

%%%%%%%%%%%%%%%%%%%%%%%%%%
%Definicion de la portada%
%%%%%%%%%%%%%%%%%%%%%%%%%%
\begin{titlepage}
    \begin{center}
	\begin{tabular}{ccc}
	    \includegraphics[width=3cm]{img/utfsm}
	    & 
	    \hspace{-0.2cm}
	    \begin{tabular}{c}
		Universidad T\'ecnica Federico Santa Mar\'ia \\ \hline
		\vspace{0.2cm}
		Departamento de Inform\'atica\\
		\vspace{1.2cm}
	    \end{tabular}
	    \hspace{0.2cm}
	    &
            \includegraphics[width=2cm]{img/di}
	\end{tabular}

	\vspace{1cm}
	%Titulo del Documento
	\begin{tabular}{c}
		\Huge{\sc{Testing, Detection and Possible}}\\\\
		\Huge{\sc{Solutions for the Bufferbloat}}\\\\
		\Huge{\sc{Phenomenon on Networks.}}\\\\
		\Large{\sc{State of Art}}\\\\
	\end{tabular}

    	\vspace{2cm}
	\Large{\sc{JUAN S. CATALAN OLMOS}}\\
	\vspace{2cm}
	\large{Definici\'on de Tema de Memoria}\\
	\large{para optar al T\'itulo de:}\\
	\large{INGENIERO CIVIL INFORMATICO}\\
	\vspace{2cm}
         	\normalsize{Profesor Gu\'ia: Horst H. von Brand Skopnik}\\
         	\normalsize{Profesor Coreferente: Ra\'ul Monge Anwandter}\\
	%Fecha
    
    \vspace{2cm}
    \normalsize{\sc{\today}}\\
    \normalsize{\sc{Valpara\'iso, Chile}}\\
    \end{center}
\end{titlepage}


\newpage

\section{Problem Definition}
If a little salt makes food taste better, then a lot must make it taste great, right. What happend if you apply the same statement to a network domain? It keeps been as good as it was? It improves the performance or makes it worse?.\\

\textit{Lets think of a network as a road system where everyone drives at the maximum speed. When the road gets full, there are only two choices: crash into other cars, or get off the road and wait until things get better. The former isn't as disastrous on a network as it would be in real life: losing packets in the middle of a communication session isn't a big deal. But making a packet wait for a short time is usually better than ``dropping" it and having to wait for a retransmission.}[iv]\\

At this point, the role of the router becomes important. It has to control the congestion effectively in networks. It is important to remember that the traffic in a network is inherently bursty, so the role of the buffers in the router is to smooth the flow of traffic. Without any buffering, to allocate the bandwidth evenly would be impossible. But there are some problems with current algorithms; they use tail-drop based queue management that has two big drawbacks: 1.- lockout 2.- full queue that impact with a high queue delay.\\

These problems are fixed with the creation of a group of FIFO based queue management mechanisms to support end-to-end congestion control in the internet. That procedure is called Active Queue Management (AQM). With AQM the loss of package and the average queue length is reduced; this impacts in a decreasing end-to-end delay by droping packages before buffer comes full, using the exponential weighted average queue length as a congestion indicator. For the proper use of AQM, it has to be widely enabled and consistently configured the router.\\

\textit{Today's networks are suffering from unnecessary latency and poor system performance. The culprit is Bufferbloat, the existence of excessively large and frequently full buffers inside the network. Large buffers have been inserted all over the Internet without sufficient thought or testing. They damage or defeat the fundamental congestion-avoidance algorithms of the Internet's most common transport protocol. Long delays from bufferbloat are frequently attributed incorrectly to network congestion, and this misinterpretation of the problem leads to the wrong solutions being proposed.}[i]\\

The existence of cheap memory and a misguided desire to avoid packet loss has led to larger and larger buffers being deployed in the hosts, routers, and switches that make up the Internet. It turns out that this is a recipe for bufferbloat. Evidence of bufferbloat has been accumulating over the past decade, but its existence has not yet become a widespread cause for concern.\\

This phenomenon lately has begun to being a concern and you can find in websites like \url{http://lwn.net}, publications as \url{http://queue.acm.org} or with a simple google' search many hits and mail lists related that threats the subject. All this let to the creation of a specific website [v], where anyone who is intrested in the topics related with TCP/IP networks but mostly with \textit{Bufferbloat} can join. All of these work, led us to the work of whose been the leader and the one who show the importance that this phenomenon really has; Jim Gettys and his personal blog[iii], subject that this thesis is based on.

\section{General Objectives }
The objective of this thesis work consists of checking the effects of the existence of the \textit{BufferBloat} phenomenon, test the impact that it has in diffetent networks like residential networks, school network or/and department network and to analyze possible solutions. To accomplish this, it's first required to address the following general objectives:
\begin{center}
\begin{itemize}
\item To explain the \textit{BufferBloat} phenomenon, and explain the impact that it could have over the lantecy and throughtput in Internet.  
\item To detect the presence by a empirical mesure of the latency and throughput in a TCP/IP based LAN. 
\item To propose possible solutions in the implementation of a network where the existence of excessively large and frequently full buffers are detected, by mesuring and modeling the effects.
\end{itemize}
\end{center}

\section{Specific Objectives}
\begin{center}
\begin{itemize}
\item To select or develop appropriate test to be able to prove the existence of \textit{Bufferbloat}
\item To test and differentiate the possible cause of the excesive latency and throughput in a TCP/IP LAN and proof how much is generated by the \textit{Bufferbloat} or by a miss-configuration. 
\item To propose a possible configuration of the TCP parameters in a Linux based machine or an algorithm that can help to minimize the phenomenon. 
\end{itemize}
\end{center}

\section{Workflow}
The following is a list with the main task to do in order to accomplish this thesis, with an estimation of period of time (in weeks) that will take to do each, and how the work is going to be distributed:
\begin{itemize}
\item Definition and understanding of the key concepts related with the \textit{BufferBloat} phenomenon.( 3 )
\item Research and develop the state of art of the \textit{BufferBloat} phenomenon and the related technologies.( 4 )
\item To develop and apply different kind of tests to detect the existence of the phenomenon. ( 5  )
\item To mount and test different TCP configurations in a linux machine and OpenWRT router. ( 4  )
\item Analysis of results and search of possible solutions. ( 5  )
\item Final review and corrections.( 3  )
\end{itemize}

The estimated time that this thesis work is going to need to accomplish all of their objectives is 24 weeks, equivalent to 6 months, beggining in \textbf{March, 2012}.

\section{References}
\begin{enumerate}[i.-]
\item \textit{Bufferbloat: Dark Buffers in the Internet}\\
	Author: Jim Gettys, Kathleen Nichols\\
	Published: Queue vol. 9, no. 11, November 29, 2011\\
	Comment: This publication, Jim Gettys presents the reason and why he begun to research about this phenomenon. Also he present some fundamentals of how this phenomenon is generated and some results of his research.
\item \textit{BufferBloat: What's Wrong with the Internet?}\\
	Published: Queue vol. 9, no. 12, December 7, 2011\\
	Comment: A discussion with Vint Cerf, Van Jacobson, Nick Weaver, and Jim Gettys about the \textit{BufferBloat} phenomenon, how Gettys got motivated by Cerf to explore his home network to find the problem of ``why it was so slow", the explanation of what this phenomenon its and some other technical information related with this topic.
\item \textit{jg's Ramblings}\\
	Author: Jim Gettys\\
	\url{http://gettys.wordpress.com/}\\
	Comment: This is Gettys' personal blog. Here he write about his intrest topic. Lately he is been ``\textit{posting}" about \textit{BufferBloat}, the way he did the first experiments, how he realize that was a problem and some other technical information about this phenomenon.
\item \textit{Understanding bufferbloat and the network buffer arms race}\\
	\url{http://arstechnica.com/tech-policy/news/2011/01/}\\
	\url{understanding-bufferbloat-and-the-network-buffer-arms-race.ars}\\
	Comment: A less technical explanation of this phenomenon, the causes and how it affect the performance of a network. It makes reference mostly to the job that Gettys has done lately.
\item \textit{BufferBloat}\\
	\url{http://www.bufferbloat.net/}\\
	Comment: This is a website where you can find project related with the performance, testing and general information related with this phenomenon. Also you can find several mailing lists, git repositories, and people here that are attempting to address the Bufferbloat issues at various levels.
\end{enumerate}

\newpage
\textbf{PS:} The reason why this thesis is written in english is becouse the author wants the acceptance of his work for the international community and makes a contribution to the current work and development of this topic.\\\\
\section*{Degloce Propuesta}
\noindent
Tiempo de Escritura: 10 minutos\\ 
\section*{Degloce Anexo}
\noindent
Tiempo estimado de B\'usqueda y Lectura: 8 horas\\
Tiempo estimado de escritura: 4 horas\\
Total: \textbf{12 horas}

\end{document}
